\documentclass[11pt]{amsart}
\usepackage{amsmath,amssymb}
\usepackage{graphicx}
\usepackage[hidelinks]{hyperref}

\newcommand{\D}{\mathbb{D}}
\newcommand{\T}{\mathbb{T}}
\newcommand{\C}{\mathbb{C}}
\newcommand{\AD}{A(\mathbb{D})}
\newcommand{\SB}{S}

\title[Non-equicontinuous norm attainment in $\AD$]{A norm-attaining disc-algebra operator without boundary equicontinuity}
\author{Agent \texttt{agent\_super\_00}}
\date{2026-06-30}

\begin{document}
\maketitle

\section*{Source and Status}

This packet concerns arXiv:2305.00859, Bala--Sarkar--Sensarma,
\emph{A Bishop-Phelps-Bollobas theorem for disc algebra}~\cite{BSS}.
After Corollary 2.2, the source paper says that the authors do not know how to
remove the equicontinuity assumption from the preceding theorem, and adds that
it is unlikely that equicontinuity is necessary for norm attainment, but that
the paper has no evidence for that claim.

\begin{center}
\includegraphics[width=0.95\textwidth]{figures/open_problem_crop.png}
\end{center}

The result below is a scoped counterexample: it answers only the
norm-attainment necessity comment. It does not remove the equicontinuity
hypothesis from the Bishop-Phelps-Bollobas approximation theorem.

\section*{Definitions}

The disc algebra $\AD$ is the Banach algebra of functions continuous on
$\overline{\D}$ and holomorphic on $\D$, equipped with the supremum norm. For a
bounded operator $T:X\to \AD$, the source paper asks for equicontinuity at a
boundary point $\zeta\in\T$ of the family
\[
  \{\widetilde{Tx}:x\in S_X\}\subset C(\overline{\D}).
\]

\section*{Proof Intuition}

The obstruction is already present in the unit ball of the disc algebra. The
monomials $z^n$ have norm one and oscillate faster and faster near every
boundary point. The identity operator therefore attains its norm at the
constant function $1$, while its image of the unit sphere contains a sequence
that forbids equicontinuity at any boundary point.

\section*{Theorem}

There is a norm-attaining operator $T:\AD\to\AD$ such that
$\{\widetilde{Tf}:f\in S_{\AD}\}$ is not equicontinuous at any point of $\T$.

\begin{proof}
Let $T=I_{\AD}$ be the identity operator on $\AD$. Then $\|T\|=1$, and $T$
attains its norm at the constant function $1$, since $\|T1\|_\infty=1$.

Fix an arbitrary boundary point $\zeta\in\T$. For each $n\geq 1$, define
\[
  f_n(z)=\overline{\zeta}^{\,n}z^n\qquad(z\in\overline{\D}).
\]
Then $f_n\in S_{\AD}$ and $f_n(\zeta)=1$. Put
\[
  z_n=\zeta e^{i\pi/n}.
\]
Then $z_n\in\T$, $z_n\to \zeta$, and
\[
  f_n(z_n)=\overline{\zeta}^{\,n}(\zeta e^{i\pi/n})^n=e^{i\pi}=-1.
\]
Thus
\[
  |f_n(z_n)-f_n(\zeta)|=2
\]
while $z_n\to\zeta$. Consequently the family $S_{\AD}=I_{\AD}(S_{\AD})$ is not
equicontinuous at $\zeta$. Since $\zeta$ was arbitrary, equicontinuity fails at
every boundary point.
\end{proof}

\section*{Verification Notes and Scope}

The proof is analytic and uses no computation beyond rendering the source crop.
The counterexample is intentionally narrow. It proves that boundary
equicontinuity is not a necessary condition for norm attainment of an operator
into $\AD$. It does not prove the full Bishop-Phelps-Bollobas conclusion for
arbitrary operators into $\AD$, nor does it show that the source theorem's
equicontinuity hypothesis can be removed.

Cheap run-index search found no prior packet for arXiv:2305.00859. Exact web
searches for the source title together with \emph{disc algebra},
\emph{equicontinuity}, \emph{norm attaining}, and
\emph{Bishop-Phelps-Bollobas} did not surface a later settlement or this exact
observation. Because the example is elementary, the main human-review question
is whether the packet's scope matches the source sentence.

\section*{Human Review Recommendation}

Review as a scoped counterexample to the source paper's norm-attainment
necessity comment. Do not count it as a solution of the stronger open problem
of removing equicontinuity from the BPB theorem.

\begin{thebibliography}{9}
\bibitem{BSS}
N. Bala, J. Sarkar, and A. Sensarma,
\emph{A Bishop-Phelps-Bollobas theorem for disc algebra},
arXiv:2305.00859.
\end{thebibliography}

\end{document}
